\scanpage{99}
% archive sheet number: 98

\uline{Paramètres locaux} \uline{du site modulaire} :

a$'$) Dans le complémentaire de $M'\cup M''\cup M^\infty$ $\left(\underset{\substack{|\\c_4=0}}{j=0},\,\underset{\substack{|\\c_6=0}}{j=12^3},\,\underset{\substack{|\\\Delta=0}}{j=\infty}\right)$ \deleted{est} \deleted{\unsure{on|la}} on peut prendre la fonction $j$

\deleted{\unsure{a)|b)}} Dans le complémentaire de $M'\cup M''$, on peut prendre $1/j$ ou $\dfrac{1}{j-12^3}$.

\deleted{b) Dans le complémentaire de}

b) Au voisinage de $(M')_6$, on peut prendre $c_4$.

c) \rule{4em}{.3pt} $(M'')_6$, \rule{4em}{.3pt} $c_6$.

d) Au voisinage du point $p=2$, $c_4=0$, $c_6=0$ on peut prendre $a_1$

e) Au voisinage du point $p=3$, $c_4=0$, $c_6=0$ \rule{4em}{.3pt} $b_2$

\uline{Question} Donner la structure de l'anneau local complété de $M$ en un point géom. où $c_4=0$, $c_6=0$ [il y en a \deleted{\unsure{essentiellement}]}\ deux : \unsure{Bien|Bon} \unsure{sûr|pris}, l'un en car. $2$, l'autre en car. $3$]
\[
\deleted{\unsure{A=|R=}}\,\underset{\substack{|\\\text{corps résiduel}\\\text{clôture algébrique}\\\text{de }\mathbb F_2\ \text{resp. }\mathbb F_3}}{W}[[t]]
\qquad (t=a_1\ \text{en car. }2,\quad t=b_2\ \text{en car. }3)
\]
et les équations des diviseurs $c_4=0$, $c_6=0$ \ldots..

\noindent\rule{\linewidth}{.3pt}

\margin{$\left.\begin{gathered}
G\ \text{d'ordre }48\\
B\ \text{d'ordre }12\\
G/B\ \text{d'ordre }4\\
B/\mu_6\ \text{d'ordre }2\simeq\mu_2\\
\mu_6=\underbrace{\mu_2}_{\substack{\Vert\\\text{centre}\\ Z}}\times\underbrace{\mu_3}_{\substack{\Vert\\\unsure{B_u|Bu}}},\\
G'=\mathrm{Sl}(2,\mathbb Z/3\mathbb Z),\ B'=G'\cap B\\
G'/B'\ \text{d'ordre }4\ldots\\
G'\ \text{d'ordre }24\\
B'\ \text{d'ordre }6
\end{gathered}\right\}$}
% slanted

Étudions la structure de $V(c_4)$ dans $(M_3)$. On trouve que si $E$ est une courbe elliptique sur $S$, avec $6$ inv. sur $S$ et $c_4(E)=0$, alors ds $\mathord{}_3E$ il y a un unique drapeau invariant par $\underline{\mathrm{Aut}}_S(E)\simeq\mu_6$. Ceci permet de ``réduire le groupe structural'' $G=\mathrm{Gl}(2,\mathbb Z/3\mathbb Z)$ de $(M_3)[1/6]$ de $G$ à $B$ (le Borel). On trouve que sur $\underline S=\mu_3^*[\frac16]$ (schéma des racines primitives 3ièmes de $1$), \deleted{\unsure{qui comprend|qui}} il y a une courbe elliptique $\mathbb E$ avec $c_4(E)=0$, et une base $e_1^0,e_2^0$ adaptée au drapeau précédent, $\langle e_1^0,e_2^0\rangle=\zeta^0$ (où $\zeta^0$ est la racine 3ième de l'unité canonique) ; cette $\overline E$ est définie à isom. unique près. Toute courbe elliptique \inserted{$E$} sur $S$, avec rigidification \inserted{$e_1,e_2$} d'échelon 3, définit un $S\longrightarrow\underline S$ (\unsure{après|grâce à} : $\zeta=\langle e_1,e_2\rangle$), \deleted{et} $E$ est isomorphe à $\overline E\times_{\underline S} S$, avec \inserted{(loc.)} soit la rigidification \unsure{image inverse}, soit \deleted{la rigidification}\inserted{une des trois autres} qui s'en déduisent par translation par élmt de $G$ \unsure{mod.|m. de|mod} $B$. Donc $M_3[1/6]\cap V(c_4)$ est isomorphe à la somme de $4$ copies de $\mu_3^*[\frac16]$, et est can. isom. à $\mu_3^*(1/6)\deleted{\unsure{\times}} \mathbin{\overset B\times}G$, comme espace où $G$ opère ($B$ opère sur $\mu_3^*$ via déterminant). T.S.V.P.
